\documentclass[12pt,a4paper]{article}
\usepackage[utf8]{inputenc}
\usepackage[spanish]{babel}
\usepackage{amsmath}
\usepackage{amsfonts}
\usepackage{amssymb}
\usepackage{graphicx}
\usepackage{booktabs}
\usepackage{geometry}
\geometry{a4paper, margin=1in}

\title{Reporte de Avance: Detección Automatizada de Tigmotaxis \\ \large Validación Inicial y Trabajo a Futuro}
\author{Trabajo Terminal}
\date{\today}

\begin{document}

\maketitle

\section{Introducción}
El presente reporte detalla la metodología y los resultados preliminares obtenidos en el desarrollo de un modelo computacional para la detección automatizada de Tigmotaxis en roedores (un indicador conductual de ansiedad). Se ha utilizado un pipeline híbrido que integra aprendizaje profundo para el seguimiento postural (DeepLabCut) y aprendizaje automático supervisado (SimBA) para la clasificación de la conducta basada en geometría y cinemática, con el objetivo de dotar a las evaluaciones neurofarmacológicas de una validez objetiva y escalable.

\section{Metodología y Procesamiento}

\subsection{Inferencia Postural con DeepLabCut}
Se implementó el modelo preentrenado \textit{SuperAnimal TopViewMouse} (basado en ResNet50) para extraer 8 puntos de interés anatómico (body-parts) de los sujetos experimentales en videos de alta resolución (720p a 30 FPS). Para abordar la enorme carga computacional inherente al análisis de registros ininterrumpidos de 5 minutos (11,100+ fotogramas), se ideó una arquitectura de script en módulos aislados ("atomic processing") para la iteración de inferencias en GPU (NVIDIA RTX 5070 Ti). De este modo, se garantizó la salud y la liberación sistemática de la memoria VRAM entre procesos, evitando latencias de fragmentación e interrupciones del kernel.

\subsection{Preprocesamiento y Mitigación de Ruido (Anti-Jitter)}
Para atenuar eficientemente las oscilaciones falsas o ``jitter'' inherentes al rastreo top-view, las coordenadas vectoriales generadas (\textit{.h5}) fueron interceptadas y procesadas a través de un esquema algorítmico que implementa filtros iterativos de \textit{Savitzky-Golay}. Esto resulta imperativo para consolidar las derivadas de la cinemática animal (gradientes de velocidad angular y varianzas de aceleración), que representan las características base de SimBA para discernir la conducta.

\subsection{Sincronización de Metadatos y Definición de Regiones de Interés (ROIs)}
La Tigmotaxis se cataloga invariablemente por la afinidad espacial del sujeto biológico por aferrarse o resguardarse sobre la zona tangencial (periferia / muros) de un entorno confinado. Durante la integración de los videos completos de 5 minutos, la plataforma SimBA presentó una validación rígida de precondiciones, invisibilizando los archivos en la sección de dibujo de ROIs debido a la ausencia de indexación previa. \textbf{Justificación algorítmica:} Para sortear este bloqueo de la interfaz sin requerir manipulación manual registro por registro, fue imperativo desarrollar un script de sincronización (\texttt{registrar\_videos\_simba.py}). Este módulo empleó la librería \textit{OpenCV} para auditar nativamente la resolución y el \textit{framerate} real, inyectando estos metadatos forzosamente en el log estructural \texttt{video\_info.csv} asumiendo la escala basal calibrada probada empíricamente (400 mm a 2.5 px/mm). Sorteada la restricción de visibilidad, se definieron meticulosamente los contornos geométricos de las arenas experimentales listos para el análisis.

\subsection{Procesamiento Analítico Automatizado (API de SimBA)}
Tradicionalmente, la extracción de métricas derivadas en SimBA exige navegar y ejecutar procesos dependientes de botones visuales a través de su entorno gráfico, lo cual resulta en un cuello de botella insostenible al procesar lotes masivos de iteraciones experimentales repetitivas interdiarias. \textbf{Justificación algorítmica:} Con la consigna final de automatizar permanentemente el \textit{pipeline} y erradicar este factor obstructivo; se proyectó una herramienta puente vía terminal de comandos (CLI) denominada \texttt{automatizar\_simba.py}.  Este orquestador enlazó directamente con la base programática (backend) en Python de SimBA instanciando las clases abstractas maestras \texttt{ExtractFeaturesFrom8bps} y \texttt{ROIFeatureCreator}. Esta abstracción abstrae al analista humano posibilitando el procesamiento masivo escalable e ininterrumpible de archivos \textit{.csv} hacia las topologías espaciales que exige finalmente el motor de Random Forest de manera desatendida.

\section{Evaluación Preliminar (Prueba Ciega Diagnóstica)}
Pudiendo prever las deficiencias de inducción estadística de un modelo piloto originado sobre secuencias fragmentarias de 40 segundos, consideramos prudente implementar una audición con un video incógnito y masivo de 5 minutos, correspondiente a un set experimental con suministro de un ansiolítico benzodiacepínico (\texttt{DZP-R1.mov}).

Utilizando rutinas implementadas en OpenCV de manera pragmática, se compuso una secuencia probatoria visual integrando un Heads-Up Display (HUD) con cronometraje acumulativo y termometría probabilística de la Tigmotaxis cuadro a cuadro. \textbf{Justificación algorítmica:} Dado que el Random Forest inicial había sido alimentado con un acervo conductual minúsculo carente de sesgos reales inter-etológicos en transiciones largas de tiempo libre; arrastraba latente el riesgo de censurar pasivamente micromovimientos no contundentes por falsos negativos (FN).  Para potenciar el diagnóstico y erradicar la miopía del clasificador, se efectuó metodológicamente sobre el script \texttt{predict\_proba} el descenso artificial del umbral intrínseco de \textit{certeza paramétrica de corte} fijado visualmente en $0.18$ (18\% de tolerancia de escrutinio), estimulando artificialmente su reactividad para visibilizar y registrar imperfecciones que se debatirán clínicamente de forma ocular o auditiva comparando a la predicción emitida por la Inteligencia Artificial. Posteriormente, se evaluaron rigurosamente los resultados producidos con base a este enfoque.

\begin{table}[h!]
\centering
\begin{tabular}{lc}
\toprule
\textbf{Métrica Diagnóstica Observada} & \textbf{Conteo} \\
\midrule
Verdaderos Positivos (Eventos asertivos) & 8 \\
Falsos Positivos & 14 \\
Falsos Negativos (Ceguera del modelo a la conducta) & 6 \\
Micro-Detecciones con escisión artificial (``Flickering'' temporal $\sim$0.2s) & 4 \\
\bottomrule
\end{tabular}
\caption{Resultados del Análisis Cualitativo en video ciego (\texttt{DZP-R1.mov} de 5 min).}
\end{table}

\subsection{Diagnóstico del Comportamiento del Modelo}
A partir del análisis expuesto, la capacidad inferencial basal del modelo exhibe una valía incuestionable para interpretar posturas orientadas a la Tigmotaxis. Sin embargo, su calibración evidencia severidades metodológicas que restringen su escalabilidad:
\begin{enumerate}
    \item La propensión alta a la asunción de \textit{Falsos Positivos} sugiere una sobresensibilidad o "falacia algorítmica", atribuible a un entrenamiento desproporcionado comparado a su dieta deficiente de ejemplos del comportamiento base ininterrumpido a lo largo del encuadre temporal real de 5 minutos completos en arena de experimentación. El modelo clasifica de manera azarosa a través del ruido del movimiento natural animal y a su escasa biblioteca analítica.
    \item La activación inestable, el quiebre y las fisuras de \textit{Micro-detecciones} exponen la precarización predictiva cuadro-por-cuadro, producto de la nulidad de un dictamen regulado por intervalos mínimos de conducta temporal persistente y contigua (sin constreñimientos para dictaminar conductas minúsculas atípicas de 200 ms).
\end{enumerate}

\section{Plan de Acción y Trabajo a Futuro}
Las correcciones y las implementaciones para la siguiente iteración experimental de entrenamiento persiguen directamente extirpar la precarización diagnóstica expuesta:

\begin{itemize}
    \item \textbf{Sustitución Sistemática a Entrenamiento de Amplio Espectro (5 Minutos):} Se dejarán progresivamente en obsolescencia los clips recortados para entrenar en 5 videogramas ininterrumpidos ya registrados en las rutinas masivas (\texttt{C1-R1}, \texttt{C2-R1}, \texttt{C56-R1}, \texttt{C7-R1}, \texttt{R5DZ}).
    \item \textbf{Human-in-The-Loop:} Empleando conocimiento etológico experto durante la categorización aural extensa ($5 \times \sim300\text{s}$), se espera enriquecer colosalmente la heterogeneidad de ejemplos del Random Forest, precipitando sustancialmente el ratio de los falsos positivos experimentados en ambientes cerrados.
    \item \textbf{Ajustes Hiperparamétricos (Minimum Bout Length):} Sobre la consola heurística de SimBA, se habilitará la supresión paramétrica para intervalos conductuales $\leq 0.3 s$. Esta restricción diluirá orgánicamente las interrupciones del \textit{Flickering}, obligando a la regresión a suavizar su toma de variables lógicas y de continuidad clínica en lugar de juicios fraccionados por parpadeos musculares del roedor.
\end{itemize}

El cumplimiento de esta directriz promete decantar un modelo excepcionalmente curado para generar estadísmos puros de \textit{ansiedad inducida / control \textit{vs.} tratada}, logrando justificar así la automatización integral de los cribados preclínicos documentados para el proyecto.

\end{document}
